% AI-produced by GPT-6 Astra (OpenAI, ChatGPT Pro mode). No priority claim.
\subsection{Totality does not give compact convex inverse limits}

The following result rules out a possible averaging shortcut in the
AKR persistence problem and its proposed strengthening under
totality. The same positive-approximation obstruction holds in
every repaired AKR category and in the total predimension category
$\mathcal C_\delta$.

\begin{proposition}[Proved: a uniform obstruction to positive approximation]
\label{CC:positive-gap}
Let $A$ be either a repaired AKR category $K^*$ covered by
Proposition~\ref{C:akr-all-rings}, or the ordered predimension
category $\mathcal C_\delta$. For every object $X$ and every
finitely supported family of nonnegative real coefficients
$(t_g)_{g:W\to X}$ with $\sum_g t_g=1$, put
\[
 \mu_i(v)=\sum_{g:\,gv_i=v}t_g\quad
(v\in A(V,X),\ i=1,2,3).
\]
Then
\[
 \|\mu_2-\mu_1\|_1+\|\mu_3-\mu_1\|_1\geq1.
\]
In particular the initial triple has no arbitrarily accurate
simultaneous equalizers obtained from probability combinations of
ordinary arrows, even as the target varies.
\end{proposition}
\begin{proof}
For either category, every nonempty finite set $S$ of $W$-arrows
into one object satisfies
\[
 \left|\bigcup_{g\in S}\{gv_1,gv_2,gv_3\}\right|>\tfrac32|S|.
\]
For $K^*$ this is the proved counting invariant from
Proposition~\ref{C:akr-all-rings}. For $\mathcal C_\delta$, the
arrows in $S$ have distinct hyperedge images, since an increasing
three-point embedding is determined by its image. If $U$ is
their vertex union, positive predimension gives
\[
 2|U|>3|E_X(U)|\geq3|S|.
\]
Every point in this union is represented by a composite $gv_i$,
even though other raw vertices of $X$ need not be $d$-closed.
In both cases distinct underlying points give distinct point maps,
since the category is concrete with singleton $V$.
For a point map $v$, set
\[
 w(v)=\max\bigl(\{t_g:gv_i=v\text{ for some }i\}\cup\{0\}\bigr).
\]
Apply the counting invariant to the finite threshold sets
$S_s=\{g:t_g>s\}$. Integrating the elementary identities for the
lengths of the intervals $0<s<t_g$ and $0<s<w(v)$ gives
\[
 \sum_v w(v)
 =\int_0^\infty\left|\bigcup_{g\in S_s}g(W)\right|\,ds
 \geq\frac32\int_0^\infty|S_s|\,ds
 =\frac32.
\]
All sums and integrals reduce to finite step functions on the support
of $t$. Nonnegativity gives
$w(v)\leq\max_i\mu_i(v)$. Moreover,
\[
 \max_i\mu_i(v)\leq
 \mu_1(v)+(\mu_2(v)-\mu_1(v))_+
              +(\mu_3(v)-\mu_1(v))_+.
\]
Every $\mu_i$ has total mass one. Thus the sums of the positive and
negative parts of $\mu_i-\mu_1$ agree, and
\[
 \frac32\leq\sum_v\max_i\mu_i(v)
 \leq1+\frac12\bigl(
 \|\mu_2-\mu_1\|_1+\|\mu_3-\mu_1\|_1\bigr).
\]
Rearrangement proves the bound. The integrated inequality is in
fact strict for each finite probability support; only its uniform
lower bound is used. No positive lower bound on
$\inf_m a_m-3/2$ is required.
\end{proof}

\begin{theorem}[Refuted: totality guarantees compact convex limits]
\label{CC:convex-counterexample}
There is a nonempty connected quasi-filtered small category $A$ and
a contravariant diagram of nonempty finite-dimensional compact convex
sets, with continuous affine transition maps, whose strict inverse
limit is empty. One may take $A=K^*$ as above, or the countable
category $A=\mathcal C_\delta$, which is total by
Theorem~\ref{A:predimension-total}. Thus the proposed implication
fails even under totality.
\end{theorem}
\begin{proof}
For either choice of $A$, write $E_X=A(W,X)$ and $P_X=A(V,X)$.
Both sets are finite because the objects are finite sets and all
arrows are injective functions. In the finite-dimensional real vector
space $\mathbb R[P_X]^2$ put
\[
 D_X[g]=([gv_2]-[gv_1],\,[gv_3]-[gv_1])
 \qquad(g\in E_X).
\]
If $E_X$ is nonempty, Proposition~\ref{CC:positive-gap} says that the
compact convex hull $C_X=\operatorname{conv}\{D_X[g]:g\in E_X\}$
does not meet the open unit ball for the $\ell^1$ norm.
Finite-dimensional convex separation gives a nonzero linear functional
$\lambda$ with
\[
 \inf_{z\in C_X}\lambda(z)
 \geq\sup_{\|b\|_1<1}\lambda(b)=\|\lambda\|_\infty.
\]
Here one separates a convex set from a disjoint open convex set; the
last equality is the elementary duality of the two specified norms.
Dividing by $\|\lambda\|_\infty$ gives coefficient functions
$\phi_1,\phi_2:P_X\to[-1,1]$ such that
\begin{equation}\label{CC:potential-inequality}
 \phi_1(gv_2)-\phi_1(gv_1)
 +\phi_2(gv_3)-\phi_2(gv_1)\geq1
 \qquad(g\in E_X).
\end{equation}

Define $\mathcal C(X)$ to be the set of all pairs
$(\phi_1,\phi_2)\in[-1,1]^{2P_X}$ satisfying
\eqref{CC:potential-inequality}. It is a closed convex subset of a
finite-dimensional cube and is nonempty by the preceding separation.
If $E_X$ is empty, the conditions are vacuous, so use the whole cube,
which is nonempty also when $P_X$ is empty.

For an arrow $u:X\to Y$, define
\[
 \mathcal C(u):\mathcal C(Y)\longrightarrow\mathcal C(X),
 \qquad(\phi_i)_i\longmapsto
 \bigl(v\longmapsto\phi_i(uv)\bigr)_i.
\]
This is continuous and affine. It preserves the inequalities because
the inequality for $g:W\to X$ after restriction is the inequality
for $ug:W\to Y$ before restriction. Restrictions compose in the
required order, so $\mathcal C$ is the asserted contravariant diagram.

Suppose a compatible family $(\phi_{i,X})_{X,i}$ existed. For each
$v:V\to X$, compatibility along $v$ evaluated at $1_V$ would give
\[
 \phi_{i,X}(v)=\phi_{i,V}(1_V).
\]
Hence each coefficient function is constant on all point maps, with
the same constant across objects. At $X=W$, the inequality for
$g=1_W$ would then read $0\geq1$, a contradiction. Thus the limit
is empty. Quasi-filteredness and connected nonemptiness of $K^*$ are
the established properties of the repaired AKR construction;
the corresponding properties and totality of $\mathcal C_\delta$
are proved in Theorem~\ref{A:predimension-total}.
\end{proof}

\paragraph{Scope for the persistence problem.}
The theorem refutes the general compact convex inverse-limit lemma
suggested as a way to extend locally bounded cycle functionals,
including its proposed totality-specific version.
A particular diagram of such functionals could still have a compatible
point for additional reasons. Neither quasi-filteredness nor totality
alone supplies the required compatible bounded potentials. The
construction does not assume the
stronger compact Hausdorff condition of
Theorem~\ref{CD:compact-Ramsey}.
